% Extracted from p04-depth-width
\subsection{Depth--Width Tradeoffs in Looped Architectures}
\label{ch-depthwidth}

\paragraph{Background}

The Universal Transformer \citep{dehghani2019universal} replaced
fixed depth with recurrent iteration of a shared block. ALBERT
\citep{lan2020albert} showed cross-layer parameter sharing
preserves quality while shrinking models. The OpenMythos family
\citep{gomez2024openmythos} combines recurrence with sparse
experts in a prelude/loop/coda structure. Test-time loop
extrapolation---running extra loops at inference---is the natural
test-time scaling axis for such architectures, though
recent small-scale evidence finds gains saturating at the
training loop count (Section~\ref{sec:loopsweep}).

\paragraph{The budget equations}
\label{sec:budget}

Let $L_u$ be the number of unique non-looped layers (prelude $+$
coda), $L_b$ the unique layers in the looped block, $R$ the loop
count, $d$ the width, and $p \approx 12d^2$ the parameters per
dense layer (attention $4d^2$ plus feed-forward $8d^2$; MoE
multiplies the feed-forward term by the expert count). Splitting
off fixed costs (embeddings, output head, norms) as $P_0$,
\begin{equation}
  P = P_0 + (L_u + L_b)\,p, \qquad
  D_{\mathrm{eff}} = L_u + R\,L_b,
  \label{eq:budget}
\end{equation}
where $D_{\mathrm{eff}}$ is the \emph{effective} (unrolled) depth:
the number of layer evaluations per forward pass. Eliminating
$L_b$,
\begin{equation}
  D_{\mathrm{eff}} = R\,\frac{P - P_0}{p} - (R - 1)\,L_u .
  \label{eq:deff}
\end{equation}
Two readings of \eqref{eq:deff}: at fixed $R$, effective depth
grows with the parameter budget as usual; at fixed $P$,
effective depth grows \emph{linearly in $R$}. Loops manufacture
depth out of thin parameter air.

The catch is compute. Forward FLOPs per token scale as
$2\,p\,D_{\mathrm{eff}}$ (plus small terms), and training as
roughly $6\,p\,D_{\mathrm{eff}}$. A dense model with $P$
parameters spends $\approx 6P$ FLOPs per token; a looped model
with the same $P$ spends $6\,p\,D_{\mathrm{eff}} > 6P$, because
shared parameters are reused. \emph{Loops do not save compute;
they convert compute into depth at fixed parameter count.} Any
Chinchilla-style compute-optimality analysis must be re-derived
for looped models, since parameters and FLOPs no longer move
together.

\paragraph{The conditioning constraint}
\label{sec:conditioning}

Depth bought with loops has a price beyond FLOPs: gradient flow.
Backpropagation through $R$ repetitions of the loop block
multiplies by its Jacobian $J$ at each pass, so gradient norms
scale roughly as $\rho^R$, where $\rho$ is the spectral radius of
$J$. If $\rho > 1$, gradients explode with $R$; if $\rho < 1$
they vanish and the later loops learn nothing. Stable, useful
looping requires $\rho \approx 1$---the block must act as a
near-isometry that still does computational work each iteration.

This yields a structural argument for \emph{small} loop blocks. A
one- or two-layer operator is a tighter dynamical system: fewer
interacting nonlinearities, easier to keep near-isometric, and a
cleaner fixed-point-iteration semantics (refine the
representation, don't rebuild it). A deep looped block must
simultaneously learn a complex transformation \emph{and} keep its
$R$-fold composition well-conditioned---a harder joint problem.
Hence our first hypothesis:

\paragraph{Hypothesis H1.} At fixed $P$ and fixed training
compute, configurations with small $L_b$ and moderate $R$
outperform large $L_b$ with $R = 1$, because the conditioning
constraint binds before the expressivity constraint.

The prelude and coda play different roles from the loop block:
the prelude maps tokens into the latent space the loop refines,
and the coda maps back to logits. Neither is iterated, so
neither faces the $\rho^R$ constraint, but starving them
bottlenecks everything the loops can do. Hence:

\paragraph{Hypothesis H2.} There is a minimal viable $L_u$ (we
conjecture 4--8 layers at $d = 512$ scale); below it, loop-count
increases cannot compensate, because the prelude/coda mappings
themselves lack capacity.

\paragraph{Width, loops, and the compute knob}
\label{sec:width}

With $L_u$ floored by H2 and $L_b$ kept small by H1, the remaining
budget goes to width $d$ and loop count $R$. Width buys
representational capacity per step; loops buy sequential
refinement steps. They are not interchangeable: width helps
single-pass pattern matching, loops help iterative computation.
This suggests treating $R$ not as an architectural constant but
as a \emph{compute knob} akin to test-time compute
\citep{snell2024scaling}:

\paragraph{Hypothesis H3.} The optimal $R$ scales with task
difficulty (reasoning depth required), not with model size. A
model trained with randomised loop counts should therefore
extrapolate to $R$ beyond its mean training value, whereas
fixed-$R$ training calibrates the dynamics to exactly $R$---which
is consistent with loop-sweep measurements showing perplexity
saturating precisely at the training loop count (29.38 at 1 loop
$\to$ 4.41 at 4 loops, then flat) on our 44M model.

H3 reframes the negative extrapolation result as expected
behaviour rather than a failure: fixed-$R$ training teaches the
block a 4-step operator, and there is no reason a 4-step operator
iterated 8 times should improve.

\paragraph{Proposed experiments}

We propose, but do not run: (E1) a fixed-$P$, fixed-compute sweep
over $(L_u, L_b, R)$ on TinyStories \citep{eldan2023tinystories},
measuring validation loss, to test H1/H2 directly; (E2)
tracking the spectral radius $\rho$ of the loop-block Jacobian
through training, testing whether stable runs indeed hold
$\rho \approx 1$ and whether instability predicts divergence;
(E3) training with $R$ sampled per batch and sweeping $R$ at
evaluation, testing H3's extrapolation claim.

\paragraph{Limitations}

Entirely theoretical: no training runs are reported and the
hypotheses are untested. The $\rho^R$ argument treats the loop
block Jacobian as fixed, ignoring its evolution during training
and input dependence. TinyStories exercises limited reasoning
depth, so H3's difficulty-scaling claim needs harder benchmarks
to bite.

\paragraph{Conclusion}

Parameter sharing splits the depth question in two: unique
layers, bounded by the budget equation \eqref{eq:budget}, and
effective depth, purchasable with loops at the price of compute
and conditioning. The analysis points to small loop blocks, a
minimal but non-trivial prelude/coda, width for capacity, and
loop count as a difficulty-matched compute knob. Whether the
hypotheses survive contact with experiment is now an empirical
question---the equations say exactly which experiments to run.
